\documentclass[11pt]{article}
\usepackage[margin=1in]{geometry}
\usepackage{graphicx}
\usepackage{amsmath}
\usepackage{booktabs}
\title{A Finite Intercept Correction Calculation}
\author{Anonymous}
\date{}
\begin{document}
\maketitle
\begin{abstract}
% Fill the abstract.
\end{abstract}
\section{Introduction}

Examining residuals provides a direct way to describe discrepancies between observed targets and model predictions. Foundational treatments develop numerical and graphical residual analysis and extend the definition of residuals beyond linear models \cite{Anscombe1963ExaminationAnalysis,Cox1968GeneralDefinition}. Broader regression diagnostics address influential observations, model specification, and the interpretation of discrepancies \cite{Belsley1980RegressionDiagnostics,Bollen1985RegressionDiagnostics,Davison1992RegressionModel}. These perspectives support a modest teaching objective: use a transparent calculation to show precisely which discrepancy an additive correction removes. Such an objective is consistent with the emphasis on parsimonious models and attention to model inadequacies in statistical practice \cite{Box1976ScienceStatistics}.

The interpretation of a prediction error depends on both the discrepancy measure and the observations being assessed. Work on model validation and comparison discusses difference measures and the limitations of relying on correlation alone \cite{Willmott1981ValidationModels,Willmott1985StatisticsEvaluation}. Forecast verification frameworks distinguish attributes of predictions and observations, while mean-square-error skill decompositions explicitly include systematic and nonsystematic bias terms \cite{Murphy1987GeneralFramework,Murphy1988SkillScores}. Here, however, a systematic offset means only a constant residual on a specified finite sample. It should not be confused with the statistical bias of an estimator across repeated samples, a different concept discussed in the bias--variance literature \cite{Geman1992NeuralNetworks}.

This note considers exactly five inputs,
\[
x=(-2,-1,0,1,2),
\qquad y_i=2x_i+3,
\]
and a supplied fixed predictor $\hat y_i^{(0)}=2x_i$. Every observed residual is therefore $r_i=y_i-\hat y_i^{(0)}=3$. The calculation takes their mean,
\[
b=\frac{1}{5}\sum_{i=1}^{5}(y_i-2x_i)=3,
\]
and adds this scalar to the unchanged predictor. On these observations, the corrected prediction $2x_i+b$ equals the target at every point. The experimental log records mean squared errors of $9$, $2.25$, and $0$ at $b=0$, $1.5$, and $3$, respectively; the corresponding mean absolute errors are $3$, $1.5$, and $0$. These values describe one deterministic calculation implemented in base R.

The central interpretive constraint is that the same five targets determine and evaluate the correction. Apparent-error analysis explains why using observations both to fit a rule and assess its accuracy can yield optimistic assessments of future prediction error \cite{Efron1986HowBiased}. Cross-validatory assessment and predictive sample reuse address different evaluation questions from reporting a fitted-sample discrepancy \cite{Stone1974CrossvalidatoryChoice,Geisser1975PredictiveSample}. Subsequent treatments connect prediction-error estimation to covariance penalties and review the properties of cross-validation procedures \cite{Efron2004EstimationPrediction,Arlot2010SurveyCrossvalidation}. No such prediction-error estimation procedure is performed here. The recorded zero error establishes agreement on the supplied observations only.

The purpose of this teaching note is to make the arithmetic and its evidential scope explicit: a constant observed residual is removed by adding its mean, and the resulting error is assessed on the observations used to calculate that mean. The only recorded comparison is among the three supplied values of $b$. The calculation establishes neither a literature gap nor an algorithmic advance, and it provides no held-out evidence or comparative performance claim against the contextual methods discussed below.

\section{Related Work}

\paragraph{Intercept adjustment with a fixed predictor.}
The choice to include an intercept has a substantial methodological and educational literature. Treatments of regression without an intercept discuss its applicability, fitting considerations, and evaluative statistics \cite{Hahn1977FittingRegression,Hawkins1980NoteFitting,Eisenhauer2003RegressionThrough}. Analyses of regression software and alternative definitions of $R^2$ identify difficulties that arise when interpreting models fitted through the origin \cite{Gordon1981ErrorsComputer,Kvlseth1985CautionaryNote}. A leverage-based interpretation also connects regression through the origin to fitting an augmented dataset \cite{Casella1983LeverageRegression}. These works supply context for distinguishing intercept inclusion from intercept omission. The present calculation is narrower: the slope is supplied as $2$, and only an additive scalar is determined from the observed residuals.

Broader regression frameworks provide useful distinctions between changing a constant and changing a predictive function. Generalized linear models accommodate systematic effects through suitable transformations, generalized additive models replace linear covariate effects with estimated smooth functions, and gradient boosting develops stagewise additive expansions through optimization in function space \cite{Nelder1972GeneralizedLinear,Hastie1986GeneralizedAdditive,Friedman2001GreedyFunction}. Ridge regression and data-augmentation approaches concern biased estimation and prediction in multiple regression \cite{Hoerl1970RidgeRegression,Allen1974RelationshipBetween}. None of these broader fitting procedures is implemented here. For the supplied residuals, the relevant algebra is simply
\[
\frac{1}{5}\sum_{i=1}^{5}(r_i-b)^2
=
\frac{1}{5}\sum_{i=1}^{5}(r_i-\bar r)^2
+(b-\bar r)^2,
\qquad
\bar r=\frac{1}{5}\sum_{i=1}^{5}r_i.
\]
Thus, choosing $b=\bar r$ removes the mean-residual component of squared error on this sample. It removes every residual here because all five residuals equal $3$.

Residual magnitude should also be distinguished from influence and leverage. Work on influential observations, the hat matrix, and relationships among regression diagnostics examines how observations affect fitted quantities \cite{Cook1977DetectionInfluential,Hoaglin1978HatMatrix,Chatterjee1986InfluentialObservations}. Posterior predictive assessment instead evaluates model fitness using discrepancies and model-based reference distributions \cite{Gelman1996PosteriorPredictive}. The present note performs neither an influence analysis nor a probabilistic model check. Its constant residual vector is directly available from the supplied targets and predictions.

Forecasting provides a more application-specific setting for intercept corrections. Research on economic forecasting discusses sources of forecast uncertainty and intercept adjustment, including the effects of structural change and possible trade-offs between forecast bias and variance \cite{Clements1994TowardsTheory,Clements1996InterceptCorrections,Clements1998ForecastingEconomic}. Judgmental adjustments in macroeconomic forecasting and statistical relationships between numerical model outputs and weather observations illustrate other forms of forecast modification \cite{Turner1990RoleJudgement,Glahn1972UseModel}. These approaches address forecasting settings and assumptions absent from the five-point example. The example contains no temporal observations or structural-change experiment, so its exact correction supplies no evidence about robustness in those settings.

Prediction-model validation and updating offer another relevant perspective. Logistic-model validation, prognostic-model assessment, performance frameworks, and calibration guidance distinguish several aspects of predictive adequacy \cite{Miller1991ValidationTechniques,Harrell1996MultivariablePrognostic,Steyerberg2009AssessingPerformance,Initiative2019CalibrationAchilles}. Work on survival-model calibration, logistic-model updating, selection among update strategies, and reviews of statistical updating methods describes alternatives ranging from limited recalibration to broader model revision \cite{Houwelingen2000ValidationCalibration,Steyerberg2004ValidationUpdating,Vergouwe2016ClosedTesting,Su2016ReviewStatistical}. Intercept updating is therefore a contextual connection, but the present response-scale addition is not an implementation of these clinical prediction procedures. It uses no probability link, survival model, or independent validation population.

Post-processing also appears in probability and uncertainty calibration. Sigmoid mappings of support-vector-machine scores, comparisons of probability-calibration methods, and temperature scaling concern the relationship between classifier outputs and outcome probabilities \cite{Platt1999ProbabilisticOutputs,Niculescumizil2005PredictingGood,Guo2017CalibrationModern}. Probability-score partitioning separates uncertainty, reliability, and resolution \cite{Murphy1973NewVector}. Regression uncertainty calibration and conformal predictive inference address uncertainty estimates or prediction bands rather than merely correcting a point-prediction offset \cite{Kuleshov2018AccurateUncertainties,Levi2022EvaluatingCalibrating,Lei2018DistributionfreePredictive}. These are distinct objectives. Zero point-prediction error on the supplied sample establishes no probability-calibration or predictive-coverage result.

\paragraph{Residual-derived correction and evaluation.}
The distinction between fitting and prediction is central to interpreting the calculation. Regression shrinkage research explicitly considers deterioration from original-data fit to new-data prediction, while work on prediction-rule error estimation and discriminant-analysis error rates studies how assessment procedures behave \cite{Copas1983RegressionPrediction,Efron1983EstimatingError,Lachenbruch1968EstimationError}. Information criteria account for adjusted parameters in model identification, and generalized degrees of freedom quantify the sensitivity of fitted values to observed data \cite{Akaike1974NewLook,Ye1998MeasuringCorrecting}. These perspectives motivate making the dependence of $b$ on the targets explicit. They do not provide an estimated optimism correction or a future-error estimate for this deterministic example.

Studies of regression submodel selection and comparisons of cross-validation with bootstrap assessment investigate evaluation beyond apparent fit \cite{Breiman1992SubmodelSelection,Kohavi1995CrossvalidationBootstrap}. Research on using cross-validation for model selection shows that tuning and evaluation can interact, and later analysis clarifies which prediction-error quantities cross-validation estimates and the difficulties of associated uncertainty intervals \cite{Varma2006BiasError,Bates2023CrossvalidationWhat}. A survey of leakage in machine-learning-based science further documents the consequences of inappropriate information dependencies for reported conclusions \cite{Kapoor2023LeakageReproducibility}. Here, the dependency is deliberately disclosed: observed residuals determine $b$, and those observations are reused for assessment. No cross-validation, bootstrap assessment, or held-out experiment was run, so these references contextualize the scope of the reported errors rather than validate generalization.

Finally, discussions comparing mean absolute error with root mean square error emphasize that error summaries have different interpretations \cite{Willmott2005AdvantagesMean}. This note reports the supplied mean absolute and mean squared errors without claiming that one metric is universally preferable. Both vanish at $b=3$ because each of the five corrected residuals is zero. The literature supplies conceptual context; the experimental evidence remains exclusively the recorded finite calculation.
\section{Method}
% Fill the method.
\section{Experiments}
% Fill the experiments.
\section{Conclusion}
% Fill the conclusion.
\bibliographystyle{plain}
\bibliography{references}
\end{document}